% Pursuing Stacks (cotes 134-2 to 134-8): archive cote 134-6, site page 60 (PDF page 61), typescript page 378.
% typescript. Draft compared: ps5.tex:4381-4460.
% AI-transcribed by Claude Opus 5.5 (claude-opus-5-5, adaptive thinking, effort high), Message Batches API, batch ps3 (msgbatch_01XdFmuVVBAcgc7EFxdtCqsj), 2026-11-12; not yet checked by eye.
\scanpage{134-6/60}
\mspage{378}

102.    We now focus attention \over{\unsure{on}}{upon} the pair of mutually dual \inserted{small} categories
\[ (1) \qquad (A,B) \ , \quad \text{with} \ \ B = A^{o} \ \ \text{i.e.} \ \ A = B^{o} \ \ , \]
and recall the equivalence \inserted{of section 93,} following from the universal property of Add(B)\hand{\unsure{,}}
\[ \xout{A^{\wedge}_{ab}xxxx\underline{Hom}(Add(B),x(Ab)x\underline{Hom}(Add(B),x(Ab)x)xx} \]
\[ (2) \qquad A^{\wedge}_{ab} = \underline{Hom}(A^{o},(Ab)) \over{\simeq}{\approxeq} \underline{Homadd}(Add(\over{B}{A^{o}}), \ (Ab)) \quad , \]
which we parallel with the formula (3) of section \deleted{\unsure{§}}99 (p. 362),  which
reads when replacing in it B by A
\[ A^{\wedge}_{ab} \simeq \underline{Homadd}( \ Add(A)^{o}, \ (Ab)) \quad ; \]
\xout{xxxxx}\inserted[typed]{this} immediately suggests a canonical equivalence of categories
\[ \xout{(2)} \qquad \xout{\underline{Hom}Add(A^{o}x)xxxx\underline{Homadd}(A)x} \]
\[ (\xout{2}3) \qquad Add(A^{o}) \ \underline{\approx} \ Add(A)^{o} \quad , \]
\deleted{and similarly}
\[ \deleted{(3') \qquad Addinf(A^{o}) \simeq Addinf(A)^{o} \ ,} \]
\deleted{these equivalences} following immediately \inserted{indeed} from the \xout{xx} 2-universal proper-%
ties of these categories.   We complement (2) by the similar formula
\[ (4) \qquad F \mapsto \widetilde{F} : \qquad A^{\wedge}_{ab} \xrightarrow{\approx} \underline{Homaddinf}(Addinf(B), \ (Ab)) \quad , \quad B = A^{o} \ , \]
where \uline{Homaddinf} denotes the category of infinitely additive functors
between two \xout{addi} infinitely additive categories.  In view of the emphasis
lately on chain complexes in Addinf(B) rather than in Add(B),  in order
to reconstruct say the derived category $D.(B^{\wedge}_{ab})$ of chain complexes in
$B^{\wedge}_{ab}$ ,  and get existence of ``integrators'' with \xout{values}\inserted[typed]{components} in Addinf(B)
(whereas there may be none with \xout{values} components in Add(B)),  it is
f\unsure{or}mula (4) rather than (2) which is going to be relevant for our homolo-%
gy formalism.  Using (4),  we get \xout{for any object L in Addinf(B)}\inserted[typed]{a canonical biadditive pairing}
\[ (*) \qquad A^{\wedge}_{ab} \ x \ Addinf(B) \longrightarrow (Ab) \ , \]
which visibly is exact with respect to the first \deleted{left hand} factor, and
which we may equally interprete as a functor
\[ (5) \qquad L \mapsto \widetilde{L} \ \ ; \qquad Addinf(B) \longrightarrow \deleted{\unsure{\underline{Homadd}}} \ \ \underline{Homex}(A^{\wedge}_{ab},(Ab)) \quad , \]
where \uline{Homex}  denotes the category of \uline{exact} (hence additive)  functors \xout{of}
from an abelian category to another one.\margin{//} It can be shown that the pai-%
ring (*)  can be extended canonically to a pairing
\[ (6) \qquad A^{\wedge}_{ab} \ x \ B^{\wedge}_{ab} \longrightarrow (Ab) \]
commuting to small direct limits in each variable,  and identifying (up

\margin{60}
